\textbf{Phần 1. Trắc nghiệm lựa chọn}
\setcounterpageref{ex}{0}

\begin{ex} %[5H8N46]
$1\text{ dm}^3$ bằng bao nhiêu xăng-ti-mét khối?
\choice
{\True $1\,000\text{ cm}^3$}
{$10\text{ cm}^3$}
{$100\text{ cm}^3$}
{$10\,000\text{ cm}^3$}
\loigiai{
Ta có: $1\text{ dm}^3 = 1\,000\text{ cm}^3$.\\
\textbf{Đáp án của câu hỏi là A}
}
\end{ex}

\begin{ex} %[5H8N46]
$1\text{ dm}^3$ bằng bao nhiêu lít?
\choice
{$10\text{ l}$}
{$0{,}1\text{ l}$}
{\True $1\text{ l}$}
{$100\text{ l}$}
\loigiai{
Ta có: $1\text{ dm}^3 = 1\text{ l}$.\\
\textbf{Đáp án của câu hỏi là C}
}
\end{ex}

\begin{ex} %[5H8N46]
Số đo “Bốn nghìn hai trăm năm mươi xăng-ti-mét khối” được viết là
\choice
{$425\text{ cm}^3$}
{$4\,250\text{ dm}^3$}
{$42\,500\text{ cm}^3$}
{\True $4\,250\text{ cm}^3$}
\loigiai{
Viết là $4\,250\text{ cm}^3$.\\
\textbf{Đáp án của câu hỏi là D}
}
\end{ex}

\begin{ex} %[5H8H46]
Đổi $3{,}5\text{ dm}^3$ ra xăng-ti-mét khối ta được
\choice
{$35\text{ cm}^3$}
{$350\text{ cm}^3$}
{\True $3\,500\text{ cm}^3$}
{$35\,000\text{ cm}^3$}
\loigiai{
$3{,}5\text{ dm}^3 = 3{,}5 \times 1\,000 = 3\,500\text{ cm}^3$.\\
\textbf{Đáp án của câu hỏi là C}
}
\end{ex}

\begin{ex} %[5H8H46]
Đổi $7\,800\text{ cm}^3$ ra đề-xi-mét khối ta được
\choice
{$0{,}78\text{ dm}^3$}
{\True $7{,}8\text{ dm}^3$}
{$78\text{ dm}^3$}
{$780\text{ dm}^3$}
\loigiai{
$7\,800\text{ cm}^3 = 7\,800 : 1\,000 = 7{,}8\text{ dm}^3$.\\
\textbf{Đáp án của câu hỏi là B}
}
\end{ex}

\begin{ex} %[5H8H46]
Kết quả của phép tính $450\text{ cm}^3 + 2{,}5\text{ dm}^3$ là
\choice
{\True $2\,950\text{ cm}^3$}
{$452{,}5\text{ cm}^3$}
{$700\text{ cm}^3$}
{$295\text{ cm}^3$}
\loigiai{
Đổi $2{,}5\text{ dm}^3 = 2\,500\text{ cm}^3$.\\
Phép tính: $450 + 2\,500 = 2\,950\text{ (cm}^3\text{)}$.\\
\textbf{Đáp án của câu hỏi là A}
}
\end{ex}

\begin{ex} %[5H8H46]
Một chai nước ngọt có dung tích $1{,}5\text{ l}$. Dung tích chai nước đó theo xăng-ti-mét khối là
\choice
{$150\text{ cm}^3$}
{$15\,000\text{ cm}^3$}
{$15\text{ cm}^3$}
{\True $1\,500\text{ cm}^3$}
\loigiai{
$1{,}5\text{ l} = 1{,}5\text{ dm}^3 = 1\,500\text{ cm}^3$.\\
\textbf{Đáp án của câu hỏi là D}
}
\end{ex}

\begin{ex} %[5H8V46]
Tính giá trị biểu thức: $12\text{ dm}^3 \times 4 - 8\,500\text{ cm}^3$.
\choice
{$40\text{ dm}^3$}
{$38{,}5\text{ dm}^3$}
{\True $39{,}5\text{ dm}^3$}
{$395\text{ dm}^3$}
\loigiai{
$12\text{ dm}^3 \times 4 = 48\text{ dm}^3$.\\
$8\,500\text{ cm}^3 = 8{,}5\text{ dm}^3$.\\
$48 - 8{,}5 = 39{,}5\text{ dm}^3$.\\
\textbf{Đáp án của câu hỏi là C}
}
\end{ex}

\begin{ex} %[5H8V46]
Một bể cá chứa $45\text{ dm}^3$ nước, người ta đổ thêm vào bể $15\,000\text{ cm}^3$ nước. Số lít nước có trong bể lúc này là
\choice
{\True $60\text{ l}$}
{$50\text{ l}$}
{$55\text{ l}$}
{$65\text{ l}$}
\loigiai{
$15\,000\text{ cm}^3 = 15\text{ dm}^3 = 15\text{ l}$.\\
Tổng số nước: $45 + 15 = 60\text{ l}$.\\
\textbf{Đáp án của câu hỏi là A}
}
\end{ex}

\begin{ex} %[5H8V46]
Điền dấu thích hợp vào chỗ chấm: $5\text{ dm}^3\; 45\text{ cm}^3 \;\dots\; 5{,}45\text{ dm}^3$.
\choice
{$>$}
{\True $<$}
{$=$}
{Không so sánh được}
\loigiai{
$5\text{ dm}^3\; 45\text{ cm}^3 = 5{,}045\text{ dm}^3 < 5{,}45\text{ dm}^3$.\\
\textbf{Đáp án của câu hỏi là B}
}
\end{ex}

\begin{ex} %[5H8V46]
Một hình hộp chữ nhật có thể tích $3{,}6\text{ dm}^3$, diện tích đáy là $180\text{ cm}^2$. Chiều cao của hình hộp chữ nhật là
\choice
{$2\text{ cm}$}
{$200\text{ cm}$}
{$0{,}2\text{ cm}$}
{\True $20\text{ cm}$}
\loigiai{
Đổi $3{,}6\text{ dm}^3 = 3\,600\text{ cm}^3$.\\
Chiều cao là: $3\,600 : 180 = 20\text{ (cm)}$.\\
\textbf{Đáp án của câu hỏi là D}
}
\end{ex}

\begin{ex} %[5H8C46]
Một bình tràn chứa đầy nước. Người ta thả chìm $4$ viên bi sắt có thể tích như nhau vào bình thì thấy tràn ra $0{,}32\text{ dm}^3$ nước. Thể tích của mỗi viên bi sắt là
\choice
{$8\text{ cm}^3$}
{\True $80\text{ cm}^3$}
{$800\text{ cm}^3$}
{$0{,}08\text{ cm}^3$}
\loigiai{
Đổi $0{,}32\text{ dm}^3 = 320\text{ cm}^3$.\\
Thể tích mỗi viên bi sắt: $320 : 4 = 80\text{ (cm}^3\text{)}$.\\
\textbf{Đáp án của câu hỏi là B}
}
\end{ex}

\textbf{Phần 2. Trắc nghiệm đúng sai}
\setcounterpageref{ex}{0}

\begin{ex} %[5H8H46]
Xét tính đúng sai của các phép đổi đơn vị sau:
\choiceTF
{\True $4\text{ dm}^3 = 4\,000\text{ cm}^3$.}
{$2{,}05\text{ dm}^3 = 2\,500\text{ cm}^3$.}
{\True $350\text{ cm}^3 = 0{,}35\text{ dm}^3$.}
{\True $8\text{ dm}^3\; 6\text{ cm}^3 = 8{,}006\text{ dm}^3$.}
\loigiai{
\begin{itemize}
    \item $4 \times 1000 = 4000\text{ cm}^3$. Đúng.
    \item $2{,}05 \times 1000 = 2050\text{ cm}^3 \neq 2500\text{ cm}^3$. Sai.
    \item $350 : 1000 = 0{,}35\text{ dm}^3$. Đúng.
    \item $8 + 6 : 1000 = 8{,}006\text{ dm}^3$. Đúng.
\end{itemize}
\begin{itemize}
    \item a) \textbf{Đúng}.
    \item b) \textbf{Sai}.
    \item c) \textbf{Đúng}.
    \item d) \textbf{Đúng}.
\end{itemize}
}
\end{ex}

\begin{ex} %[5H8V46]
Một xô nước chứa $12\text{ l}$ nước. Người ta dùng một ca có dung tích $400\text{ cm}^3$ để múc nước từ xô ra. Xét tính đúng sai:
\choiceTF
{\True Dung tích xô nước bằng $12\text{ dm}^3 = 12\,000\text{ cm}^3$.}
{\True Dung tích chiếc ca là $0{,}4\text{ l}$.}
{Cần múc $25$ ca đầy thì hết nước trong xô.}
{\True Cần múc $30$ ca đầy thì hết nước trong xô.}
\loigiai{
\begin{itemize}
    \item $12\text{ l} = 12\text{ dm}^3 = 12\,000\text{ cm}^3$. Đúng.
    \item $400\text{ cm}^3 = 0{,}4\text{ dm}^3 = 0{,}4\text{ l}$. Đúng.
    \item Số ca cần múc: $12\,000 : 400 = 30$ ca $\neq 25$ ca. Ý 3 Sai, Ý 4 Đúng.
\end{itemize}
\begin{itemize}
    \item a) \textbf{Đúng}.
    \item b) \textbf{Đúng}.
    \item c) \textbf{Sai}.
    \item d) \textbf{Đúng}.
\end{itemize}
}
\end{ex}

\textbf{Phần 3. Trả lời ngắn}
\setcounterpageref{ex}{0}

\begin{ex} %[5H8H46]
Đổi $6{,}75\text{ dm}^3$ ra xăng-ti-mét khối.
\shortans{6750}
\loigiai{
\textbf{Đáp số:} $6750$\\
$6{,}75 \times 1\,000 = 6\,750\text{ cm}^3$.
}
\end{ex}

\begin{ex} %[5H8H46]
Đổi $12\,400\text{ cm}^3$ ra đề-xi-mét khối (dùng dấu phẩy nếu có).
\shortans{12,4}
\loigiai{
\textbf{Đáp số:} $12{,}4$\\
$12\,400 : 1\,000 = 12{,}4\text{ dm}^3$.
}
\end{ex}

\begin{ex} %[5H8V46]
Tính kết quả phép tính: $3{,}2\text{ dm}^3 + 4\,800\text{ cm}^3$ theo đơn vị đề-xi-mét khối.
\shortans{8}
\loigiai{
\textbf{Đáp số:} $8$\\
$3{,}2 + 4{,}8 = 8\text{ dm}^3$.
}
\end{ex}

\begin{ex} %[5H8V46]
Một hình hộp chữ nhật có thể tích là $5{,}4\text{ dm}^3$. Biết diện tích đáy là $300\text{ cm}^2$. Tính chiều cao hình hộp chữ nhật theo xăng-ti-mét.
\shortans{18}
\loigiai{
\textbf{Đáp số:} $18$\\
$5{,}4\text{ dm}^3 = 5\,400\text{ cm}^3 \Rightarrow 5\,400 : 300 = 18\text{ cm}$.
}
\end{ex}

\begin{ex} %[5H8V46]
Một bình nước đang chứa $1{,}8\text{ dm}^3$ nước. Người ta rót đều vào $9$ chiếc cốc như nhau. Hỏi mỗi cốc chứa bao nhiêu xăng-ti-mét khối nước?
\shortans{200}
\loigiai{
\textbf{Đáp số:} $200$\\
$1{,}8\text{ dm}^3 = 1\,800\text{ cm}^3 \Rightarrow 1\,800 : 9 = 200\text{ cm}^3$.
}
\end{ex}

\begin{ex} %[5H8C46]
Một bể cá mini hình hộp chữ nhật có diện tích đáy $250\text{ cm}^2$. Thả một khối đá vào bể thì nước dâng lên thêm $1{,}6\text{ cm}$ (khối đá ngập hoàn toàn). Tính thể tích khối đá theo đề-xi-mét khối.
\shortans{0,4}
\loigiai{
\textbf{Đáp số:} $0{,}4$\\
Thể tích khối đá: $250 \times 1{,}6 = 400\text{ cm}^3 = 0{,}4\text{ dm}^3$.
}
\end{ex}
